\par\indent\SourceLeaf{241}%
Combinatio huius aequationis cum aequatione (10) producit
\begingroup\fontsize{10}{12}\selectfont
\[
\begin{aligned}
DD''-D'D'&=(\alpha\alpha''+\beta\beta''+\gamma\gamma''-\alpha'\alpha'-\beta'\beta'-\gamma'\gamma')\Delta\\
&\quad+E(n'n'-nn'')+F(nm''-2m'n'+mn'')+G(m'm'-mm'')
\end{aligned}
\]
\endgroup
Iam patet esse
\[
\begin{gathered}
\frac{\mathrm{d}E}{\mathrm{d}p}=2m,\quad\frac{\mathrm{d}E}{\mathrm{d}q}=2m',\quad\frac{\mathrm{d}F}{\mathrm{d}p}=m'+n,\\
\frac{\mathrm{d}F}{\mathrm{d}q}=m''+n',\quad\frac{\mathrm{d}G}{\mathrm{d}p}=2n',\quad\frac{\mathrm{d}G}{\mathrm{d}q}=2n''
\end{gathered}
\]
sive
\[
\begin{array}{lll}
m=\tfrac{1}{2}\dfrac{\mathrm{d}E}{\mathrm{d}p},&m'=\tfrac{1}{2}\dfrac{\mathrm{d}E}{\mathrm{d}q},&m''=\dfrac{\mathrm{d}F}{\mathrm{d}q}-\tfrac{1}{2}\dfrac{\mathrm{d}G}{\mathrm{d}p}\\[6pt]
n=\dfrac{\mathrm{d}F}{\mathrm{d}p}-\tfrac{1}{2}\dfrac{\mathrm{d}E}{\mathrm{d}q},&n'=\tfrac{1}{2}\dfrac{\mathrm{d}G}{\mathrm{d}p},&n''=\tfrac{1}{2}\dfrac{\mathrm{d}G}{\mathrm{d}q}
\end{array}
\]
Porro facile confirmatur, haberi
\[
\begin{aligned}
&\alpha\alpha''+\beta\beta''+\gamma\gamma''-\alpha'\alpha'-\beta'\beta'-\gamma'\gamma'\\
&\quad=\frac{\mathrm{d}n}{\mathrm{d}q}-\frac{\mathrm{d}n'}{\mathrm{d}p}=\frac{\mathrm{d}m''}{\mathrm{d}p}-\frac{\mathrm{d}m'}{\mathrm{d}q}\\
&\quad=-\tfrac{1}{2}\mathord{.}\frac{\mathrm{d}\mathrm{d}E}{\mathrm{d}q^2}+\frac{\mathrm{d}\mathrm{d}F}{\mathrm{d}p\mathord{.}\mathrm{d}q}-\tfrac{1}{2}\mathord{.}\frac{\mathrm{d}\mathrm{d}G}{\mathrm{d}p^2}
\end{aligned}
\]
Quodsi iam has expressiones diversas in formula pro mensura curvaturae in fine art. praec. eruta substituimus, pervenimus ad formulam sequentem, e solis quantitatibus \(E,F,G\) atque earum quotientibus differentialibus primi et secundi ordinis concinnatam:
\[
\begin{aligned}
4(EG-FF)^2k&=E(\tfrac{\mathrm{d}E}{\mathrm{d}q}\mathord{.}\tfrac{\mathrm{d}G}{\mathrm{d}q}-2\tfrac{\mathrm{d}F}{\mathrm{d}p}\mathord{.}\tfrac{\mathrm{d}G}{\mathrm{d}q}+(\tfrac{\mathrm{d}G}{\mathrm{d}p})^2)\\
&\quad+F(\tfrac{\mathrm{d}E}{\mathrm{d}p}\mathord{.}\tfrac{\mathrm{d}G}{\mathrm{d}q}-\tfrac{\mathrm{d}E}{\mathrm{d}q}\mathord{.}\tfrac{\mathrm{d}G}{\mathrm{d}p}-2\tfrac{\mathrm{d}E}{\mathrm{d}q}\mathord{.}\tfrac{\mathrm{d}F}{\mathrm{d}q}\\
&\qquad+4\tfrac{\mathrm{d}F}{\mathrm{d}p}\mathord{.}\tfrac{\mathrm{d}F}{\mathrm{d}q}-2\tfrac{\mathrm{d}F}{\mathrm{d}p}\mathord{.}\tfrac{\mathrm{d}G}{\mathrm{d}p})\\
&\quad+G(\tfrac{\mathrm{d}E}{\mathrm{d}p}\mathord{.}\tfrac{\mathrm{d}G}{\mathrm{d}p}-2\mathord{.}\tfrac{\mathrm{d}E}{\mathrm{d}p}\mathord{.}\tfrac{\mathrm{d}F}{\mathrm{d}q}+(\tfrac{\mathrm{d}E}{\mathrm{d}q})^2)\\
&\quad-2(EG-FF)(\tfrac{\mathrm{d}\mathrm{d}E}{\mathrm{d}q^2}-2\tfrac{\mathrm{d}\mathrm{d}F}{\mathrm{d}p\mathord{.}\mathrm{d}q}+\tfrac{\mathrm{d}\mathrm{d}G}{\mathrm{d}p^2})
\end{aligned}
\]
\Article{12}
Quum indefinite habeatur
\[
\mathrm{d}x^2+\mathrm{d}y^2+\mathrm{d}z^2=E\mathrm{d}p^2+2F\mathrm{d}p\mathord{.}\mathrm{d}q+G\mathrm{d}q^2
\]
patet, \(\BareRoot{(E\mathrm{d}p^2+2F\mathrm{d}p\mathord{.}\mathrm{d}q+G\mathrm{d}q^2)}\) esse expressionem generalem elementi linearis in superficie curva. Docet itaque analysis in art. praec. explicata, ad inveniendam mensuram curvaturae haud opus esse formulis finitis, quae coordina%
